% !TeX root = ../thesis.tex
% ==============================================
% CAD绘制综合参数参考附录
% —— 冷媒膨胀罐 V-1810 机械设计
% 整合设计参数、标准规范、标准件选型、关键尺寸
% ==============================================

\chapter{CAD绘制综合参数参考}

本附录为冷媒膨胀罐 V-1810 CAD工程图绘制提供集中、系统化的参数参考。囊括设计条件总览、主体材料性能、标准件选型明细、主要零部件几何参数及全部开孔方位数据，供制图成员按图索骥，避免分散查阅多本手册。

\section{设计条件总览}
\label{sec:cad-design-conditions}

\begin{table}[htbp]
  \centering
  \caption{冷媒膨胀罐 V-1810 设计条件总览}
  \label{tab:cad-design-overview}
  \small\renewcommand{\arraystretch}{1.15}
  \setlength{\tabcolsep}{6pt}
  \begin{tabularx}{\textwidth}{@{}p{4.8cm} >{\raggedright\arraybackslash}X @{}}
    \toprule
    \textbf{参数项} & \textbf{数值与说明} \\
    \midrule
    设备位号 / 名称        & 22-V-1810 / 冷媒膨胀罐 \\
    所属装置                & $100$ kt/a 釜式法 $+$ $200$ kt/a 管式法 LDPE/EVA 联合装置（18 单元） \\
    业主单位                & 江苏斯尔邦石化有限公司 \\
    建设地点                & 江苏省连云港市徐圩新区石化产业基地 \\
    压力容器类别            & 第 Ⅰ 类（TSG 21—2016，第二组介质，$PV=4.4$ MPa$\cdot$m³） \\
    操作介质                & 冷却水（冷媒），无毒、不易燃 \\
    设计使用年限            & 15 年 \\
    设备布置方式            & 露天，立式，裙座支承 \\
    \midrule
    设计压力（内压）        & $0.80$ MPa(g) \\
    设计压力（外压）        & 全真空（FV, $-0.1013$ MPa(g)） \\
    工作压力（正常操作）    & $0.60$ MPa(g) \\
    设计温度                & $170$ ℃ \\
    工作温度（正常操作）    & $155$ ℃ \\
    工作温度（温水充液）    & $100$ ℃ \\
    全真空工况金属温度      & $100$ ℃ \\
    最低设计金属温度（MDMT）& $-15$ ℃ @ $0.85$ MPa(g) \\
    \midrule
    罐体内径 $D_i$          & $1800$ mm \\
    筒体切线长度 $L$（T/T） & $1700$ mm \\
    有效容积 $V$            & $5.5$ m³ \\
    最大介质密度 $\rho$     & $1000$ kg/m³（满液校核工况） \\
    设备总高度（近似）      & $\approx 5000$ mm（含裙座） \\
    支承基础顶面标高 TOS    & $\approx 9050$ mm \\
    \midrule
    基本风压 $q_0$          & $450$ N/m²（连云港，50 年一遇，B 类地面粗糙度） \\
    抗震设防烈度            & 7 度（$0.10g$），第一组，场地类别 Ⅱ 类 \\
    极端最高气温            & $38.5$ ℃ \\
    极端最低气温            & $-12.3$ ℃ \\
    \bottomrule
  \end{tabularx}
\end{table}

\section{主体材料性能参数}
\label{sec:cad-material}

\begin{table}[htbp]
  \centering
  \caption{主体材料 Q345R 钢板（$3\sim16$ mm，正火态）性能参数}
  \label{tab:cad-q345r-full}
  \small\renewcommand{\arraystretch}{1.1}
  \setlength{\tabcolsep}{6pt}
  \begin{tabularx}{\textwidth}{@{}l l l @{}}
    \toprule
    \textbf{参数项} & \textbf{符号} & \textbf{数值} \\
    \midrule
    \multicolumn{3}{c}{\textit{力学性能（GB/T 713.1—2023）}} \\
    \midrule
    屈服强度（室温）           & $R_{\text{eL}}$   & $345$ MPa \\
    抗拉强度                   & $R_m$             & $510\sim640$ MPa \\
    断后伸长率                 & $A$               & $\ge 21\%$ \\
    $0$ ℃ 冲击吸收能量        & $KV_2$            & $\ge 41$ J \\
    弯曲试验（$180^\circ$, $d=3a$） & —             & 合格 \\
    \midrule
    \multicolumn{3}{c}{\textit{化学成分（熔炼分析，质量分数 / \%）}} \\
    \midrule
    C                          & —                 & $\le 0.20$ \\
    Si                         & —                 & $\le 0.55$ \\
    Mn                         & —                 & $1.20\sim1.70$ \\
    P                          & —                 & $\le 0.025$ \\
    S                          & —                 & $\le 0.010$ \\
    Nb                         & —                 & $0.015\sim0.050$ \\
    Alt（全铝）                & —                 & $\ge 0.020$ \\
    碳当量 $C_{\text{eq}}$    & —                 & $\le 0.45$ \\
    冷裂纹敏感性 $P_{\text{cm}}$ & —               & $\le 0.28$ \\
    \midrule
    \multicolumn{3}{c}{\textit{许用应力（GB/T 150.2—2024 附录 B）}} \\
    \midrule
    $\le 20$ ℃                & $[\sigma]$        & $189$ MPa \\
    $100$ ℃                   & $[\sigma]$        & $189$ MPa \\
    $150$ ℃                   & $[\sigma]$        & $189$ MPa \\
    $170$ ℃                   & $[\sigma]^t$      & $186$ MPa \\
    $200$ ℃                   & $[\sigma]$        & $183$ MPa \\
    \midrule
    \multicolumn{3}{c}{\textit{物理参数}} \\
    \midrule
    弹性模量（$100$ ℃）       & $E_t$             & $\approx 2.01\times10^5$ MPa \\
    密度                       & $\rho_s$          & $7850$ kg/m³ \\
    泊松比                     & $\mu$             & $\approx 0.3$ \\
    线膨胀系数（$20\sim200$ ℃）& $\alpha$          & $\approx 12.0\times10^{-6}$ /℃ \\
    钢板厚度负偏差（B类偏差）  & $C_1$             & $-0.30$ mm（固定下偏差） \\
    \midrule
    \multicolumn{3}{c}{\textit{厚度附加量}} \\
    \midrule
    腐蚀裕量（壳体/封头）      & $C_2$             & $3.0$ mm（工业冷却水，15 年） \\
    腐蚀裕量（裙座）          & $C_2$             & $2.0$ mm（大气腐蚀） \\
    \bottomrule
  \end{tabularx}
\end{table}

\begin{table}[htbp]
  \centering
  \caption{接管材料 20 号钢无缝钢管性能参数}
  \label{tab:cad-20steel}
  \small\renewcommand{\arraystretch}{1.1}
  \setlength{\tabcolsep}{8pt}
  \begin{tabularx}{\textwidth}{@{}l l l @{}}
    \toprule
    \textbf{参数项} & \textbf{符号} & \textbf{数值} \\
    \midrule
    标准                    & —        & GB 6479—2013 / GB/T 8163—2018 \\
    屈服强度（室温）        & $R_{\text{eL}}$ & $\ge 245$ MPa \\
    抗拉强度                & $R_m$           & $\ge 410$ MPa \\
    断后伸长率              & $A$             & $\ge 25\%$ \\
    许用应力（$170$ ℃）    & $[\sigma]^t_n$  & $\approx 140$ MPa \\
    许用应力（常温）        & $[\sigma]_n$    & $\approx 152$ MPa \\
    \bottomrule
  \end{tabularx}
\end{table}

\section{设计强度参数汇总}
\label{sec:cad-strength-params}

\begin{table}[htbp]
  \centering
  \caption{强度设计基础参数汇总}
  \label{tab:cad-strength-basis}
  \small\renewcommand{\arraystretch}{1.15}
  \setlength{\tabcolsep}{8pt}
  \begin{tabularx}{\textwidth}{@{}l l l l @{}}
    \toprule
    \textbf{参数} & \textbf{符号} & \textbf{数值} & \textbf{确定依据} \\
    \midrule
    设计温度许用应力                    & $[\sigma]^t$   & $186$ MPa      & GB/T 150.2—2024 附录 B，线性插值 \\
    常温许用应力                        & $[\sigma]$     & $189$ MPa      & GB/T 150.2—2024 附录 B \\
    焊接接头系数                        & $\phi$         & $0.85$         & A/B 类焊缝，局部 RT $\ge 20\%$ \\
    钢板厚度负偏差                      & $C_1$          & $0.30$ mm      & GB/T 709—2019 B类偏差（固定下偏差$-0.30$ mm） \\
    腐蚀裕量（壳体/封头）               & $C_2$          & $3.0$ mm       & HG/T 20580—2020，15 年 \\
    计算压力                            & $P_c$          & $0.80$ MPa     & 液柱静压力 $< 5\%$ 设计压力，可忽略 \\
    壳体计算厚度                        & $\delta$       & $4.57$ mm      & GB/T 150.3—2024 式 (5-1) \\
    封头计算厚度                        & $\delta_h$     & $4.56$ mm      & GB/T 150.3—2024 式 (7-1) \\
    设计厚度（壳体/封头）               & $\delta_d$     & $7.57$ mm      & $\delta + C_2$ \\
    名义厚度（壳体/封头）               & $\delta_n$     & $14$ mm        & 综合外压、刚度、开孔补强确定 \\
    有效厚度（壳体/封头）               & $\delta_e$     & $10.7$ mm      & $\delta_n - C_1 - C_2$ \\
    许用外压力                          & $[P]$          & $0.125$ MPa    & GB/T 150.3—2024 图 6-2 图算法 \\
    全真空外压载荷                      & $P_{\text{vac}}$ & $0.1013$ MPa & 1 标准大气压 \\
    外压稳定安全系数                    & $n_s$          & $1.23$         & $[P]/P_{\text{vac}}$ \\
    水压试验压力                        & $P_T$          & $1.02$ MPa     & $1.25P[\sigma]/[\sigma]^t$ 圆整 \\
    试验应力校核限值                    & —              & $263.9$ MPa    & $0.9\phi R_{\text{eL}}$ \\
    试验环向应力                        & $\sigma_T$     & $86.3$ MPa     & 合格（安全裕度 $3.06$） \\
    \bottomrule
  \end{tabularx}
\end{table}

\section{主要零部件几何参数}
\label{sec:cad-components}

\subsection{筒体}

\begin{table}[htbp]
  \centering
  \caption{筒体几何与制造参数}
  \label{tab:cad-cylinder}
  \small\renewcommand{\arraystretch}{1.15}
  \setlength{\tabcolsep}{10pt}
  \begin{tabularx}{\textwidth}{@{}l l l @{}}
    \toprule
    \textbf{参数项} & \textbf{符号} & \textbf{数值} \\
    \midrule
    材料                    & —        & Q345R 正火钢板（GB/T 713.1—2023） \\
    内径                    & $D_i$    & $1800$ mm \\
    外径                    & $D_o$    & $1828$ mm（$\delta_n = 14$ mm） \\
    名义厚度                & $\delta_n$ & $14$ mm \\
    有效厚度                & $\delta_e$ & $10.7$ mm \\
    筒体切线长度            & $L$      & $1700$ mm \\
    纵缝型式                & —        & 双面埋弧焊（SAW），全焊透对接 \\
    环缝型式                & —        & 双面埋弧焊（SAW），全焊透对接 \\
    椭圆度允差              & —        & $\le 18$ mm（$\le 1\% D_i$） \\
    棱角度允差              & —        & $\le 5$ mm \\
    成型方式                & —        & 冷卷或温卷成型 \\
    焊后热处理              & —        & SR：$600\pm20$ ℃，保温 $\ge 30$ min \\
    无损检测                & —        & 局部 RT（$\ge 20\%$，$\ge 250$ mm），Ⅲ 级合格 \\
    \bottomrule
  \end{tabularx}
\end{table}

\subsection{封头}

\begin{table}[htbp]
  \centering
  \caption{标准椭圆封头几何与制造参数}
  \label{tab:cad-head}
  \small\renewcommand{\arraystretch}{1.15}
  \setlength{\tabcolsep}{10pt}
  \begin{tabularx}{\textwidth}{@{}l l l @{}}
    \toprule
    \textbf{参数项} & \textbf{符号} & \textbf{数值} \\
    \midrule
    封头型式                & —        & 标准椭圆形封头（EHA） \\
    执行标准                & —        & GB/T 25198—2010 \\
    材料                    & —        & Q345R 正火钢板（GB/T 713.1—2023） \\
    内径                    & $D_i$    & $1800$ mm \\
    曲面深度（内高）        & $h_i$    & $450$ mm（$D_i/4$） \\
    直边段高度              & $h_s$    & $40$ mm（标准值） \\
    封头总高（含直边）      & $H_f$    & $490$ mm \\
    名义厚度                & $\delta_{nh}$ & $14$ mm（与筒体等厚） \\
    有效厚度                & $\delta_{eh}$ & $\ge 10.7$ mm（减薄后） \\
    形状系数                & $K$      & $1.0$（$D_i/2h_i = 2$） \\
    成型方式                & —        & 整体冲压或旋压成型 \\
    成型减薄率              & —        & $\approx 10\%\sim15\%$ \\
    投料厚度                & —        & $\ge 16$ mm（制造厂确定） \\
    与筒体连接              & —        & 全焊透对接接头，避开曲面过渡区 \\
    无损检测                & —        & 局部 RT，Ⅲ 级合格 \\
    \bottomrule
  \end{tabularx}
\end{table}

\subsection{裙座}

\begin{table}[htbp]
  \centering
  \caption{裙座几何与结构参数}
  \label{tab:cad-skirt}
  \small\renewcommand{\arraystretch}{1.15}
  \setlength{\tabcolsep}{10pt}
  \begin{tabularx}{\textwidth}{@{}l l l @{}}
    \toprule
    \textbf{参数项} & \textbf{符号} & \textbf{数值} \\
    \midrule
    设计标准                & —        & NB/T 47065.5—2018，参照 NB/T 47041—2014 \\
    型式                    & —        & 圆筒形焊制裙式支座 \\
    材料                    & —        & Q345R \\
    裙座筒节外径            & $D_{so}$ & $1828$ mm（与筒体外径一致） \\
    裙座筒节长度            & $H_s$    & $\approx 2000$ mm \\
    裙座壁厚（初选）        & $\delta_{ss}$ & $14$ mm \\
    腐蚀裕量                & $C_{2s}$ & $2.0$ mm \\
    与下封头连接            & —        & 全焊透对接焊接接头，位于封头直边外侧 \\
    \midrule
    \multicolumn{3}{c}{\textit{附属结构}} \\
    基础环板                & —        & 外径 $> D_{so}$，将载荷均匀传递至混凝土基础 \\
    地脚螺栓座              & —        & 径向筋板式，沿基础环板圆周均布 \\
    TOS（基础顶面标高）     & —        & $\approx 9050$ mm \\
    通气孔                  & —        & $\ge \phi 80$ mm 或等面积长圆孔，裙座上部 \\
    排液孔                  & —        & 裙座底部 \\
    检修孔                  & —        & $\ge \phi 450$ mm，裙座侧壁 \\
    涂层                    & —        & 外表面防腐涂层 + 防火涂层 \\
    \bottomrule
  \end{tabularx}
\end{table}

\section{接管与法兰标准件选型}
\label{sec:cad-nozzles}

\begin{landscape}
\begin{table}[htbp]
  \centering
  \caption{接管与法兰标准件选型明细表}
  \label{tab:cad-nozzle-flange}
  \small\renewcommand{\arraystretch}{1.3}
  \setlength{\tabcolsep}{3.5pt}
  \begin{tabularx}{\linewidth}{@{}c c c c c c c c c c@{}}
    \toprule
    \textbf{编号} & \textbf{用途} & \textbf{位置} & \textbf{公称直径 DN / mm} & \textbf{接管外径 $\times$ 壁厚 / mm} & \textbf{法兰类型} & \textbf{密封面} & \textbf{PN} & \textbf{法兰标准} & \textbf{接管标准} \\
    \midrule
    N1  & 冷媒进口    & 罐顶     & 100 & $\phi 114.3 \times 6.0$  & WN & RF & 20 & HG/T 20592 & GB 6479 \\
    N2  & 氮气补气口  & 罐顶     & 50  & $\phi 60.3 \times 5.0$   & WN & RF & 20 & HG/T 20592 & GB 6479 \\
    N3  & 放空口      & 罐顶     & 50  & $\phi 60.3 \times 5.0$   & WN & RF & 20 & HG/T 20592 & GB 6479 \\
    N4  & 压力表口    & 罐顶     & 25  & $\phi 33.7 \times 4.0$   & WN & RF & 50 & HG/T 20592 & GB 6479 \\
    N5  & 温度计口    & 罐顶     & 40  & $\phi 48.3 \times 5.0$   & WN & RF & 50 & HG/T 20592 & GB 6479 \\
    N6a & 液位计口（上）& 筒体   & 25  & $\phi 33.7 \times 4.0$   & WN & RF & 50 & HG/T 20592 & GB 6479 \\
    N6b & 液位计口（下）& 筒体   & 25  & $\phi 33.7 \times 4.0$   & WN & RF & 50 & HG/T 20592 & GB 6479 \\
    N7  & 压力变送器口 & 筒体    & 25  & $\phi 33.7 \times 4.0$   & WN & RF & 50 & HG/T 20592 & GB 6479 \\
    N8  & 冷媒出口    & 罐底     & 100 & $\phi 114.3 \times 6.0$  & WN & RF & 20 & HG/T 20592 & GB 6479 \\
    N9  & 排净口      & 罐底     & 50  & $\phi 60.3 \times 5.0$   & WN & RF & 20 & HG/T 20592 & GB 6479 \\
    M1  & 人孔        & 筒体侧壁 & 500 & $\phi 508.0 \times 12.0$ & WN & RF & 20 & HG/T 21514 & GB 6479 \\
    \bottomrule
  \end{tabularx}
\end{table}
\end{landscape}

\subsection{法兰紧固件与垫片}

\begin{table}[htbp]
  \centering
  \caption{法兰紧固件与垫片选用}
  \label{tab:cad-bolts-gaskets}
  \small\renewcommand{\arraystretch}{1.2}
  \setlength{\tabcolsep}{8pt}
  \begin{tabularx}{\textwidth}{@{}l l l @{}}
    \toprule
    \textbf{选用项} & \textbf{规格} & \textbf{标准依据} \\
    \midrule
    螺栓/螺柱型式      & 等长双头螺柱                    & HG/T 20613—2009 \\
    螺柱材质           & 35CrMoA（调质）                 & HG/T 20613—2009 \\
    螺母型式           & 六角螺母                        & HG/T 20613—2009 \\
    螺母材质           & 30CrMoA（调质）                 & HG/T 20613—2009 \\
    垫片型式           & 柔性石墨复合垫 / PTFE 包覆垫     & HG/T 20606—2009 / HG/T 20607—2009 \\
    垫片适用温度       & $-200\sim650$ ℃（柔性石墨复合垫）& — \\
    垫片适用介质       & 水、蒸汽、油品、非强氧化性介质  & — \\
    人孔密封垫片       & 柔性石墨复合垫，DN500 PN20       & HG/T 21514—2014 \\
    \bottomrule
  \end{tabularx}
\end{table}

\section{人孔标准件选型}
\label{sec:cad-manhole}

\begin{table}[htbp]
  \centering
  \caption{人孔（M1）标准件选型参数}
  \label{tab:cad-manhole-detail}
  \small\renewcommand{\arraystretch}{1.2}
  \setlength{\tabcolsep}{8pt}
  \begin{tabularx}{\textwidth}{@{}l l l @{}}
    \toprule
    \textbf{参数项} & \textbf{选型值} & \textbf{标准依据} \\
    \midrule
    人孔类型                & 回转盖带颈对焊法兰人孔     & HG/T 21514—2014 \\
    公称直径 DN             & $500$ mm                    & — \\
    公称压力 PN             & $20$（Class 150）           & — \\
    接管规格                & $\phi 508 \times 12$ mm     & — \\
    接管材料                & 20 号钢（GB 6479—2013）   & — \\
    法兰型式                & 带颈对焊法兰（WN）          & HG/T 20592—2009 \\
    法兰密封面              & 突面（RF）                  & — \\
    法兰标准号              & HG/T 20592 WN500-20 RF     & — \\
    法兰盖                  & 回转盖式，碳钢              & HG/T 21514—2014 \\
    紧固件                  & 等长双头螺柱，35CrMoA       & HG/T 20613—2009 \\
    垫片                    & 柔性石墨复合垫              & HG/T 20606—2009 \\
    补强方式                & 厚壁接管自补强              & GB/T 150.3—2024 \\
    \bottomrule
  \end{tabularx}
\end{table}

\section{开孔补强参数汇总}
\label{sec:cad-reinforcement}

\begin{table}[htbp]
  \centering
  \caption{开孔补强计算汇总}
  \label{tab:cad-reinforcement-detail}
  \small\renewcommand{\arraystretch}{1.2}
  \setlength{\tabcolsep}{3pt}
  \begin{tabularx}{\textwidth}{@{}c c c c c c c c@{}}
    \toprule
    \textbf{序号} & \textbf{接管编号} & \textbf{开孔位置} & \textbf{接管规格 DN $\times$ $\delta_t$ / mm} & \textbf{开孔直径 $d$ / mm} & \textbf{需补面积 $A$ / mm²} & \textbf{实补面积 $A_e$ / mm²} & \textbf{补强方式} \\
    \midrule
    1  & M1    & 筒体   & $500 \times 12$   & 484   & 2212  & 3647  & 厚壁接管自补强（合格） \\
    2  & N1    & 上封头 & $100 \times 6$    & 102.3 & 468   & 980   & 厚壁接管自补强（合格） \\
    3  & N8    & 下封头 & $100 \times 6$    & 102.3 & 468   & 1020  & 厚壁接管自补强（合格） \\
    4  & N2    & 上封头 & $50 \times 5$     & 50.3  & 230   & 465   & 不需另行补强（合格） \\
    5  & N3    & 上封头 & $50 \times 5$     & 50.3  & 230   & 465   & 不需另行补强（合格） \\
    6  & N9    & 下封头 & $50 \times 5$     & 50.3  & 230   & 465   & 不需另行补强（合格） \\
    7  & N4    & 上封头 & $25 \times 4$     & 25.7  & 117   & 380   & 不需另行补强（合格） \\
    8  & N5    & 上封头 & $40 \times 5$     & 38.3  & 175   & 420   & 不需另行补强（合格） \\
    9  & N6a   & 筒体   & $25 \times 4$     & 25.7  & 117   & 510   & 不需另行补强（合格） \\
    10 & N6b   & 筒体   & $25 \times 4$     & 25.7  & 117   & 510   & 不需另行补强（合格） \\
    11 & N7    & 筒体   & $25 \times 4$     & 25.7  & 117   & 510   & 不需另行补强（合格） \\
    \bottomrule
  \end{tabularx}
\end{table}

\section{接管方位与定位尺寸}
\label{sec:cad-orientation}

\begin{table}[htbp]
  \centering
  \caption{接管方位及定位尺寸（以正北为 $0^\circ$，顺时针）}
  \label{tab:cad-orientation-detail}
  \small\renewcommand{\arraystretch}{1.2}
  \setlength{\tabcolsep}{3pt}
  \begin{tabularx}{\textwidth}{@{}c c c c c c c@{}}
    \toprule
    \textbf{编号} & \textbf{用途} & \textbf{所在部件} & \textbf{方位角 / °} & \textbf{距下封头切线高度 / mm} & \textbf{外伸长度 / mm} & \textbf{内伸长度 / mm} \\
    \midrule
    N1  & 冷媒进口            & 上封头（顶轴）  & 中心     & —                       & $200$ & $0$ \\
    N2  & 氮气补气口          & 上封头          & $0$      & —                       & $200$ & $0$ \\
    N3  & 放空口              & 上封头          & $180$    & —                       & $200$ & $0$ \\
    N4  & 压力表口            & 上封头          & $60$     & —                       & $150$ & $0$ \\
    N5  & 温度计口            & 上封头          & $300$    & —                       & $250$ & $400$ \\
    N6a & 液位计口（上）      & 筒体            & $90$     & $1400$                  & $200$ & $0$ \\
    N6b & 液位计口（下）      & 筒体            & $90$     & $300$                   & $200$ & $0$ \\
    N7  & 压力变送器口        & 筒体            & $270$    & $850$                   & $150$ & $0$ \\
    N8  & 冷媒出口            & 下封头（底轴）  & 中心     & —                       & $250$ & $0$（设防涡流挡板） \\
    N9  & 排净口              & 下封头          & $180$    & —                       & $200$ & $0$ \\
    M1  & 人孔                & 筒体            & $270$    & $1000$                  & $300$ & $0$ \\
    \bottomrule
  \end{tabularx}
\end{table}

\section{管口参数汇总——DN与PN对照}
\label{sec:cad-nozzle-dn-pn}

表~\ref{tab:cad-nozzle-summary} 汇总全部管口代号与公称直径 DN、公称压力 PN 的对应关系，供 CAD 标注和材料统计直接引用。DN 按 GB/T 1047—2019 定义；PN 按 HG/T 20615—2009 的 Class/PN 对应关系确定（Class 150 $\rightarrow$ PN20，Class 300 $\rightarrow$ PN50）。

\begin{table}[htbp]
  \centering
  \caption{管口参数全表（DN与PN汇总）}
  \label{tab:cad-nozzle-summary}
  \small\renewcommand{\arraystretch}{1.35}
  \setlength{\tabcolsep}{12pt}
  \begin{tabularx}{0.7\textwidth}{@{}c c c@{}}
    \toprule
    \textbf{管口代号} & \textbf{DN} & \textbf{PN} \\
    \midrule
    A1、A2、V1、Z1    & DN50  & PN20 \\
    B1                & DN250 & PN20 \\
    L1$\sim$L4、P1、R1 & DN50  & PN50 \\
    M1                & DN600 & PN20 \\
    T1                & DN40  & PN50 \\
    \bottomrule
  \end{tabularx}
\end{table}

\section{焊缝与无损检测要求汇总}
\label{sec:cad-welds}

\begin{table}[htbp]
  \centering
  \caption{焊接接头分类及无损检测要求}
  \label{tab:cad-weld-nde}
  \small\renewcommand{\arraystretch}{1.2}
  \setlength{\tabcolsep}{5pt}
  \begin{tabularx}{\textwidth}{@{}c c c c c c@{}}
    \toprule
    \textbf{接头类别} & \textbf{焊接位置} & \textbf{焊接方法} & \textbf{接头型式} & \textbf{NDE 方法} & \textbf{检测比例 / 合格级别} \\
    \midrule
    A 类 & 筒体纵缝                      & SAW & 全焊透对接 & RT & $\ge 20\%$ 各缝长 $\ge 250$ mm，NB/T 47013.2 Ⅲ 级 \\
    A 类 & 封头拼接焊缝                  & SAW & 全焊透对接 & RT & $\ge 20\%$ 各缝长 $\ge 250$ mm，NB/T 47013.2 Ⅲ 级 \\
    B 类 & 筒体与封头环缝                & SAW & 全焊透对接 & RT & $\ge 20\%$ 各缝长 $\ge 250$ mm，NB/T 47013.2 Ⅲ 级 \\
    B 类 & 筒体与裙座连接环缝            & SAW & 全焊透对接 & RT & $\ge 20\%$，NB/T 47013.2 Ⅲ 级 \\
    C 类 & 接管与壳体角焊缝              & SMAW/GTAW & 全焊透 & PT/MT & 100\%，NB/T 47013.5 Ⅰ 级 \\
    C 类 & 法兰与接管角焊缝              & SMAW/GTAW & 全焊透 & PT/MT & 100\%，NB/T 47013.5 Ⅰ 级 \\
    D 类 & 补强圈与壳体角焊缝            & SMAW & 角焊缝 & PT/MT & 100\%，NB/T 47013.5 Ⅰ 级 \\
    D 类 & 吊耳/临时附件焊痕             & —     & —     & PT      & 100\%，NB/T 47013.5 Ⅰ 级 \\
    \bottomrule
  \end{tabularx}
\end{table}

\section{焊后热处理制度}
\label{sec:cad-pwht}

\begin{table}[htbp]
  \centering
  \caption{焊后消除应力热处理（SR）制度}
  \label{tab:cad-pwht-detail}
  \small\renewcommand{\arraystretch}{1.15}
  \setlength{\tabcolsep}{10pt}
  \begin{tabularx}{\textwidth}{@{}l l @{}}
    \toprule
    \textbf{参数项} & \textbf{规定值} \\
    \midrule
    执行标准                & HG/T 20584—2020 \\
    热处理类型              & 消除应力热处理（SR） \\
    加热温度                & $600\pm20$ ℃ \\
    保温时间                & $\ge 30$ min（按 $1$ min/mm $\times$ 14 mm $= 14$ min，取 $\ge 30$ min） \\
    升温速率（$>300$ ℃）   & $\le 200$ ℃/h \\
    降温速率（$>300$ ℃）   & $\le 260$ ℃/h \\
    冷却方式（$<300$ ℃）   & 炉内自然冷却至室温 \\
    热处理范围              & 整体炉内热处理（筒体+封头+主焊缝） \\
    热处理后检验            & 硬度检测及表面无损检测复查 \\
    \bottomrule
  \end{tabularx}
\end{table}

\section{水压试验程序参数}
\label{sec:cad-hydrotest}

\begin{table}[htbp]
  \centering
  \caption{水压试验程序与参数}
  \label{tab:cad-hydrotest-detail}
  \small\renewcommand{\arraystretch}{1.2}
  \setlength{\tabcolsep}{8pt}
  \begin{tabularx}{\textwidth}{@{}l l l @{}}
    \toprule
    \textbf{试验步骤} & \textbf{压力 / MPa} & \textbf{操作要求} \\
    \midrule
    试验介质                & —                     & 常温洁净纯水（氯离子 $\le 25$ mg/L，GB 5749—2022） \\
    介质温度                & —                     & $\ge 5$ ℃ \\
    静置脱气                & 常压                  & 注水后静置 $\ge 60$ min \\
    \midrule
    第一级升压              & $0.20$                & 保压 $10$ min，检查初始渗漏 \\
    第二级升压              & $0.51$（$0.5P_T$）   & 保压 $15$ min，检查变形情况 \\
    第三级升压              & $0.80$（设计压力）    & 保压 $20$ min，锤击检查焊缝 \\
    第四级升压              & $1.02$（试验压力）    & 保压 $\ge 30$ min，保压期间不得补压 \\
    \midrule
    降压检查                & $0.80$（设计压力）    & 全面目视及锤击检查，重点检查焊缝及密封面 \\
    最终降压                & 常压                  & 降压速率 $\le 0.1$ MPa/min \\
    罐内处理                & —                     & 排净积水，干燥压缩空气吹扫至干燥 \\
    \midrule
    合格判据                & —                     & 无渗漏、无可见变形、无异响 \\
    \bottomrule
  \end{tabularx}
\end{table}

\section{设计遵循标准体系总表}
\label{sec:cad-standards}

\begin{table}[htbp]
  \centering
  \caption{冷媒膨胀罐 V-1810 设计遵循标准体系总表}
  \label{tab:cad-standards-full}
  \small\renewcommand{\arraystretch}{1.15}
  \setlength{\tabcolsep}{4pt}
  \begin{tabularx}{\textwidth}{@{}c l >{\raggedright\arraybackslash}X @{}}
    \toprule
    \textbf{序号} & \textbf{标准编号} & \textbf{标准名称} \\
    \midrule
    1  & GB/T 150.1$\sim$150.4—2024    & 压力容器（主设计标准） \\
    2  & TSG 21—2016                    & 固定式压力容器安全技术监察规程 \\
    3  & GB/T 713.1—2023                & 承压设备用钢板和钢带 第1部分：碳钢和低合金钢 \\
    4  & GB/T 709—2019                  & 热轧钢板和钢带的尺寸、外形、重量及允许偏差 \\
    5  & GB/T 25198—2010                & 压力容器封头 \\
    6  & GB 6479—2013                   & 高压化肥设备用无缝钢管 \\
    7  & GB/T 3274—2017                 & 碳素结构钢和低合金结构钢热轧钢板和钢带 \\
    8  & GB/T 17395—2008                & 无缝钢管尺寸、外形、重量及允许偏差 \\
    9  & NB/T 47065.1-$\sim$5—2018      & 容器支座 \\
    10 & NB/T 47020$\sim$47027—2012     & 压力容器法兰、垫片、紧固件 \\
    11 & HG/T 20592$\sim$20635—2009     & 钢制管法兰、垫片和紧固件 \\
    12 & NB/T 11025—2022                & 补强圈 \\
    13 & HG/T 21514—2014                & 钢制人孔和手孔的类型与技术条件 \\
    14 & HG/T 21574—2018                & 化工设备吊耳设计选用规范 \\
    15 & NB/T 47013.1$\sim$47013.14—2015 & 承压设备无损检测 \\
    16 & HG/T 20580—2020                & 钢制化工容器设计基础规范 \\
    17 & HG/T 20584—2020                & 钢制化工容器制造技术规范 \\
    18 & HG/T 20585—2020                & 钢制低温压力容器技术规范 \\
    19 & HG/T 20553—2011                & 化工配管用无缝及焊接钢管尺寸选用系列 \\
    20 & NB/T 47041—2014                & 塔式容器 \\
    21 & NB/T 47042—2014                & 卧式容器 \\
    22 & HG 20660—2000                  & 压力容器中化学介质毒性危害和爆炸危险程度分类 \\
    23 & GB 50009—2012                  & 建筑结构荷载规范 \\
    24 & GB 50011—2010（2016 年版）     & 建筑抗震设计规范 \\
    25 & GB 5749—2022                   & 生活饮用水卫生标准（水压试验水质） \\
    26 & 3767-CD-SE-0V1810\_IS1         & 冷媒膨胀罐 V-1810 设备机械数据表（SEI） \\
    \bottomrule
  \end{tabularx}
\end{table}

\section{CAD图纸图号清单}
\label{sec:cad-drawing-list}

\begin{table}[htbp]
  \centering
  \caption{冷媒膨胀罐 V-1810 CAD图纸图号清单}
  \label{tab:cad-drawing-numbers}
  \small\renewcommand{\arraystretch}{1.2}
  \setlength{\tabcolsep}{8pt}
  \begin{tabularx}{\textwidth}{@{}c l >{\raggedright\arraybackslash}X @{}}
    \toprule
    \textbf{序号} & \textbf{图号} & \textbf{图纸名称} \\
    \midrule
    1 & 22-V-1810-00 & 冷媒膨胀罐总装图 \\
    2 & 22-V-1810-01 & 筒体图 \\
    3 & 22-V-1810-02 & 封头图（上/下） \\
    4 & 22-V-1810-03 & 裙座及基础环板图 \\
    5 & 22-V-1810-04 & 人孔（M1）部件图 \\
    6 & 22-V-1810-05 & 接管方位及管口表 \\
    7 & 22-V-1810-06 & 冷媒出口防涡流挡板详图 \\
    8 & 22-V-1810-07 & 液位计安装详图 \\
    9 & 22-V-1810-08 & 吊耳布置图 \\
    10 & 22-V-1810-09 & 基础平面及地脚螺栓布置图 \\
    \bottomrule
  \end{tabularx}
\end{table}
